\documentclass{beamer}

\usepackage[german]{babel}        % Deutschsprachige Beschriftungen
\usepackage[utf8]{inputenc}       % Utf8 Zeichensatz
\usepackage{graphicx}             % Grafiken
\usepackage{fancyvrb}             % Für VerbatimInput




% See https://www.overleaf.com/learn/latex/Beamer#:~:text=package%20in%20Overleaf-,Reference%20guide,-Below%20is%20a
\usetheme{Warsaw}
\usecolortheme{default}

\beamertemplatenavigationsymbolsempty
\setbeamertemplate{footline}[frame number]

%Information to be included in the title page:
\title{\LaTeX{} für die Maturaarbeit}
\author{Ivo Blöchliger}
\institute{Kantonsschule am Burggraben}
\date{29. August 2023}
\logo{\includegraphics[width=1.5cm]{images/ksbg-logo-no-text.pdf}}

\begin{document}

\frame{\titlepage}

\begin{frame}
\frametitle{Übersicht}
\tableofcontents
\end{frame}

\begin{frame}
	\frametitle{Download der Software}
	\begin{block}{\LaTeX{} Distribution}
		Downalod von Mik\TeX
	\end{block}
	\begin{block}{Editor}
		Download von \TeX{}Studio
	\end{block}
	Links auf https://fginfo.ksbg.ch
\end{frame}

\section{Was ist \LaTeX{}?}
\subsection{Geschichte}
\begin{frame}
	\frametitle{\LaTeX{}?}
	\begin{block}{Aussprache}
		Leitägg (oder Leitäch)
	\end{block}
	\begin{block}{Ursprung}
		\TeX{}: Textsatzsystem für Mathematik, 
		1978, Donald Knuth \\ 
		$\quad$ Ziel: Typographisch einwandfreie Dokumente\\[2mm]
		\LaTeX{}: Sammlung von Macros
		1984, Leslie Lamport
	\end{block}
	\begin{alertblock}{\includegraphics[height=1em]{images/Laser-symbol.pdf}
		Warnung \includegraphics[height=1em]{images/Laser-symbol.pdf}}
		Worddokumente können nach längerer Verwendung von \LaTeX{} zu Augenkrebs
		führen.
	\end{alertblock}
\end{frame}



\subsection{Prinzip}
\begin{frame}
	\frametitle{Prinzip von \LaTeX}
	Beschreiben was, nicht wie.
	\begin{block}{Textdatei {\tt prinzip.tex}}
		\VerbatimInput{prinzip.tex}
	\end{block}
\end{frame}


\begin{frame}
	\frametitle{Automatisieren}
	\begin{itemize}
		\item Nummerierungen
		\item Verweise
		\item Verzeichnisse
		\item Formatierungen
	\end{itemize}
\end{frame}

\begin{frame}
	\frametitle{Nummerierungen}
	\begin{block}{$\backslash$label\{bla\} und $\backslash$autoref\{bla\}}
		$\backslash$label\{bla\} {\bf nach} nummererzeugendem Befehl:\\[3mm]
		\begin{quote}
			{\tt $\backslash$section\{Blah
			Blah\}$\backslash$label\{sec:bla\}}\\[3mm]

			{\tt In $\backslash$autoref\{sec:bla\} \ldots}
		\end{quote}
	\end{block}

	\begin{block}{Labelpräfixe}
		\begin{description}
			\item[Kapitel] {\tt ch:bla}\\
			\item[Abschnitte] {\tt sec:bla}\\
			\item[Abbildungen] {\tt fig:bla}\\
			\item[Gleichungen] {\tt eq:bla}\\
		\end{description}
	\end{block}
\end{frame}

\begin{frame}
	\frametitle{Mathematik}
	Mitternachtsformel, Gleichung \ref{eq:quadsol}:
\begin{equation}\label{eq:quadsol}
	x_{1,2} = \frac{-b \pm \sqrt{b^2-4ac}}{2a}
	\qquad b^2-4ac \geq 0,\, a \neq 0
\end{equation}

	\vspace{-6mm}
	\begin{block}{\LaTeX{} Code}
		\VerbatimInput{mathe.tex}
	\end{block}
\end{frame}

\subsection{Kompilierung}
\begin{frame}
	\frametitle{Kompilierung}
	\begin{block}{Erzeugen des Dokuments}
		\begin{enumerate}
			\item Textsatz, Nummerierung, Benötigte Verweise
			\item Evtl.\ Erzeugen des Literaturverzeichnisses
			\item Textsatz, Eintragen aller Verweise, Verzeichnisse
		\end{enumerate}
	\end{block}
	\begin{example}
		Nach dem ersten Durchgang:\\[2mm]
		\begin{quote}
			\rm
			{\large \bf Abschnitt 2.1}\\
			{Wie im Abschnitt ?? ersichtlich, \ldots}
		\end{quote}
	\end{example}
\end{frame}

\begin{frame}
	\frametitle{Automatisches Kompilieren}
	\begin{itemize}
		\item Mehrere Programme nötig
		\item Richtige Reihenfolge
	\end{itemize}
	\begin{block}{Automatisieren}
		\begin{itemize}
			\item TeXStudio (Editor)
			\item Visual Studio Code Plugin (?)
			\item Makefile (make)
		\end{itemize}
	\end{block}
\end{frame}

\subsection{Bilder}
\begin{frame}
	\frametitle{Bilder}

	\begin{block}{Vektoriell bitte!}
		Wenn irgendmöglich vektorielle Bilder verwenden:
		\begin{itemize}
			\item SVG $\leftarrow$ PDF
			\item PDF (wenn nicht einfach Pixelbild)
		\end{itemize}
		Aus Office-Anwendungen: {\bf Print as PDF...}
	\end{block}
	\begin{block}{Wenn Pixel, dann}
		\begin{description}
			\item[PNG] Screenshots, Bilder mit einfarbigen Flächen
			\item[JPG] Photos
		\end{description}
	\end{block}
\end{frame}

\begin{frame}
	\frametitle{Vergleich der Formate}
	\begin{columns}[onlytextwidth]
		\begin{column}{0.33\textwidth}
		  \centering
			PDF (5.5 kB)\\[5mm]
			\includegraphics[width=\textwidth]{images/vektor.pdf}
		\end{column}
		\begin{column}{0.33\textwidth}
		  \centering
			 PNG (21 kB)\\[5mm]
			\includegraphics[width=\textwidth]{images/vektor.png}
		\end{column}
		\begin{column}{0.33\textwidth}
		  \centering
			JPG (17 kB)\\[5mm]
			\includegraphics[width=\textwidth]{images/vektor.jpg}
		\end{column}
	\end{columns}
\end{frame}

\begin{frame}
	\frametitle{Platzierung automatisch, nicht im Textfluss}
	\begin{block}{Gründe}
		\begin{itemize}
			\item Konzipiert für wissenschaftliche Publikation
			\item Abbildungen speziell farbig, eigene Seiten
			\item Layout dem Autor nicht bekannt
			\item Rücksicht auf Textfluss, Seitenfüllung
		\end{itemize}
	\end{block}
	\begin{alertblock}{Nummerieren, referenzieren!}
		\begin{itemize}
			\item Alle Abbildungen sind nummeriert,\\
		 	\item mit Legende versehen\\
			\item und {\bf im Text referenziert}.
		\end{itemize}
	\end{alertblock}
\end{frame}


\section{Software}
\subsection{Absolut nötige Software}
\begin{frame}
	\frametitle{Benötigte Software}
	\begin{itemize}
		\item \LaTeX{} Distribution (z.B. MikTeX oder TeXLive)
			\begin{itemize}
				\item Programme: {\tt pdflatex}, {\tt bibtex}, {\tt biber}, etc.
				\item \LaTeX{} Pakete: {\tt beamer}, {\tt babel}, {\tt graphicx}, etc.
			\end{itemize}
		\item Texteditor: TeXStudio, Visual Studio Code, Vim, \ldots
		\item PDF-Viewer
	\end{itemize}
\end{frame}

\subsection{Nützliche Software}
\begin{frame}
	\frametitle{Nützliche Software}
	\begin{itemize}
		\item Inkscape, LibreOffice (SVG, PDF)
		\item make o.ä. als Buildsystem
		\item Verzeichnisse/Konverter zu Bib\TeX{} Einträgen
		\item Erstellung von \LaTeX{}-Tabellen
	\end{itemize}
\end{frame}



\section{Hilfe!}
\begin{frame}
	\frametitle{Hilfe!}
	\begin{block}{Tech-Lab (E23)}
		Dienstags ab 16:40\\
		Freitags ab 17:20
	\end{block}
	\begin{block}{Teams}
		Fragen und Antworten auf dem \LaTeX{} Team
	\end{block}
	\begin{block}{Vorlage auf Github}
		Ergänzungen, Verbesserungen gerne liefern!
	\end{block}
\end{frame}

\begin{frame}
	\frametitle{Dank}
	\centering
	\Huge
	Herzlichen Dank\\
	für Ihre Aufmerksamkeit\\[1cm]
	Fragen?
\end{frame}

\begin{frame}
	\frametitle{Weiteres Vorgehen}
	\begin{block}{Plan}
		\begin{itemize}
			\item Software installieren
			\item Vorlage herunterladen
			\item Vorlage kompilieren
			\item Minimalbeispiel kompilieren
			\item Minimalbeispiel studieren/ändern
		\end{itemize}
	\end{block}
\end{frame}

\end{document}
